\section{Statutory Note: Fresh Recomputation is Intended}
\label{sec:statutory}

\subsection{The DIA's Decadal Recomputation Design}

The \dia{} specifies that the \ar{} algorithm is rerun from scratch
after each decadal census:
\begin{enumerate}
\item The new \hh{} apportionment is computed, yielding new district
      counts $k_s$ for each state.
\item The \ar{} bisection tree is constructed from the new $k_s$ values.
\item The redistricting algorithm is run with the new census data and
      the new tree.
\item The resulting map governs elections for the subsequent decade.
\end{enumerate}

This is \emph{intentional}: the DIA is not an incremental update system.
It does not try to minimise the number of tracts that change district
assignment from the previous decade.
It generates the best available map under the current census data and
current apportionment.

\subsection{Constitutional Basis}

The equal-population requirement of \textit{Wesberry v.\ Sanders}~\citep{wesberry1964}
mandates that districts be as equal in population ``as nearly as is
practicable.''
Population redistributes between census cycles: Texas grew 15\% between
2010 and 2020, while California grew 6\%.
The 2020 maps become increasingly malapportioned as the decade progresses.

The decadal fresh recomputation is the constitutional response to this
malapportionment: it resets the maps to equal population, even at the
cost of boundary disruption.
This is the system as designed, not a bug.

\subsection{The Stability Result is Reassurance, Not a Target}

The finding that scenario reapportionment disruption is far below full MCMC
decorrelation in the draft scale ($d_{\rm Ham} \leq 0.23$ vs.\
$d_{\rm Ham} \approx 1$) is a reassurance, not a design target.
The DIA does not aim to minimise boundary disruption; it aims to produce
the most compact maps consistent with the new apportionment.

The stability result says: ``As a by-product of deterministic recomputation,
the scenario can be measured and audited before it is converted into a legal
or policy claim.''
This is a feature, not a constraint.

\subsection{Election-Audit Analogy}

Election certification does not rely on a single reported total.  It reconciles
several ledgers: eligible voters, ballots issued, ballots returned, precinct
tabulations, canvass corrections, chain-of-custody records, and risk-limiting
audit samples.  The same design principle applies here.

For reapportionment stability, the audit package should reconcile:
\begin{enumerate}
\item the official apportionment count and \hh{} calculation;
\item the graph and population-input hashes;
\item the split prescription for each state and chamber;
\item the deterministic seed or seed-family manifest;
\item the old and new assignments after canonical district matching;
\item the Hamming, compactness, split, and demographic deltas; and
\item any human-facing canvass note explaining why a scenario value differs
      from the final certified value.
\end{enumerate}

This is also the natural bridge to precinct-vote auditing.  A precinct vote
audit should make its own reconciliation ledger explicit: registration
eligibility universe, ballots issued, ballots cast, rejected/provisional
ballots, tabulator totals, canvass adjustments, and sampled hand-count checks.
The transferable principle is simple: every public total should have a
lower-level manifest that lets an observer trace how the total was produced.

\subsection{Implications for DIA Drafting}

The analysis of this paper has three implications for the DIA statutory text:

\begin{enumerate}
\item \textbf{No grandfathering clause is needed.}
      Some proposals suggest that the DIA should include a ``continuity
      bonus'' rewarding plans that resemble the previous decade's map.
      This paper suggests that the \ar{} algorithm can produce
      measurable reapportionment-stability records without an explicit
      continuity bonus.
      A continuity bonus would reduce compactness without providing
      meaningful stability gains.

\item \textbf{No transition period is needed for algorithm output.}
      The new maps are immediately valid after the reapportionment.
      The 23\% Hamming distance for the Texas scenario should be treated as
      an auditable disruption estimate, not as a crisis requiring legislative
      remediation by itself.

\item \textbf{The prime factorisation changes are predictable.}
      States and administrators can pre-compute scenario tree structures
      from projection inputs years before the census.
      This allows for early public comment on the expected map changes
      before the official redistricting cycle begins.
\end{enumerate}

\subsection{A Note on the 2030 Texas Case}

Texas going from $k = 38$ to $k = 41$ (prime) is the most structurally
disruptive reapportionment in the projection set.
A binary 20/21 fallback for Texas is a qualitatively different map from
a largest-prime-first 19-container partition.
The two maps will look different to voters, even if the compactness is
similar.

This is not a failure of the algorithm; it is the correct response to the
arithmetic of apportionment.
Under the current implementation, 41 triggers the large-prime binary fallback.
There may be other policy choices for $k=41$, but they would be different
algorithms and should be named as such.
Observers should understand this before comparing the scenario to nearby
counterfactuals such as $k = 40 = 2^3 \times 5$, which would produce a
different four-level tree.
